\section{Introduction}

\subsection{Objectifs de la simmulation}

L'objectif principal de la simulation est de modéliser le problème de fourragement d'une colonie de fourmis 
artificielles sur un terrain 2D hétérogène, afin d'identifier des trajets efficaces entre la fourmilière et 
une source de nourriture. Le cadre retenu suit un modèle d'intelligence en essaim où des agents simples 
produisent un comportement collectif émergent à partir de règles locales.


Le modèle vise à reproduire une coordination indirecte entre agents via deux types de phéromones : un 
phéromone d'exploration vers la nourriture ($V_1$) et un phéromone de retour vers le nid ($V_2$). Chaque 
fourmi évolue avec deux états (\textit{unloaded}/\textit{loaded}) et applique une politique 
exploration--exploitation pilotée par le paramètre $\varepsilon$, avec des transitions d'état lors de 
l'arrivée sur la nourriture ou sur la fourmilière.


L'environnement a pour objectif de représenter une difficulté de déplacement réaliste : chaque cellule de la 
carte possède un coût de traversée dans $[0,1]$, et les zones non franchissables sont traitées comme obstacles. 
Le paysage est généré de manière fractale afin d'introduire des gradients de difficulté non uniformes et de 
contraindre la dynamique des déplacements.


Enfin, la simulation poursuit un objectif d'évaluation fonctionnelle au niveau de la colonie : la quantité 
de nourriture rapportée au nid sert d'indicateur global de performance. Ce cadre constitue la base 
expérimentale du projet, avant les étapes d'optimisation algorithmique (vectorisation et parallélisation)
 décrites dans le contexte général.